% =====================================================================
%  arquitectura.tex  --  Notas de Arquitectura y Diseno de LEIP
%  ARCHIVO UNICO: no necesita ningun otro archivo.
%  Compilar en Overleaf con pdfLaTeX (dos pasadas por el indice).
% =====================================================================
\documentclass[11pt,a4paper]{article}

% ---------------------------------------------------------------------
%  DATOS DEL PROYECTO
% ---------------------------------------------------------------------
\newcommand{\proyectoNombre}{LEIP}
\newcommand{\proyectoNombreLargo}{LatAm Economic Intelligence Platform}
\newcommand{\proyectoUniversidad}{Universidad Cat\'olica del Uruguay}
\newcommand{\proyectoFacultad}{Facultad de Ingenier\'ia y Tecnolog\'ias}
\newcommand{\proyectoCarrera}{Ingenier\'ia en Inform\'atica}
\newcommand{\docVersion}{1.2}
\newcommand{\docTitulo}{Notas de Arquitectura y Dise\~no}
\newcommand{\docFecha}{Septiembre 2026}

% ---------------------------------------------------------------------
%  PREAMBULO
% ---------------------------------------------------------------------
\usepackage[utf8]{inputenc}
\usepackage[T1]{fontenc}
\IfFileExists{lmodern.sty}{\usepackage{lmodern}\usepackage{microtype}}{}

\IfFileExists{spanish.ldf}{%
  \usepackage[spanish,es-tabla,es-noquoting]{babel}%
}{%
  \usepackage[english]{babel}%
  \renewcommand{\contentsname}{\'Indice}%
  \renewcommand{\listfigurename}{\'Indice de figuras}%
  \renewcommand{\listtablename}{\'Indice de tablas}%
  \renewcommand{\tablename}{Tabla}%
  \renewcommand{\figurename}{Figura}%
}

\usepackage[a4paper,margin=2.5cm,headheight=15pt]{geometry}
\usepackage{fancyhdr}
\usepackage{parskip}
\usepackage{booktabs}
\usepackage{tabularx}
\usepackage{longtable}
\usepackage{array}
\usepackage{enumitem}
\setlist{noitemsep,topsep=4pt}
\usepackage{graphicx}
\usepackage{xcolor}
\usepackage{tcolorbox}
\tcbuselibrary{breakable}
\usepackage[hidelinks]{hyperref}
\usepackage{url}
\urlstyle{same}

\definecolor{leipAzul}{HTML}{1F4E79}
\definecolor{leipGris}{HTML}{5A5A5A}
\definecolor{leipAmbar}{HTML}{B8860B}
\definecolor{leipAmbarBg}{HTML}{FFF8E1}

\pagestyle{fancy}
\fancyhf{}
\fancyhead[L]{\small\color{leipGris}\proyectoNombre{} --- \docTitulo}
\fancyhead[R]{\small\color{leipGris} v\docVersion}
\fancyfoot[C]{\small\thepage}
\renewcommand{\headrulewidth}{0.4pt}

\usepackage{titlesec}
\titleformat{\section}{\Large\bfseries\color{leipAzul}}{\thesection}{0.6em}{}
\titleformat{\subsection}{\large\bfseries\color{leipAzul}}{\thesubsection}{0.6em}{}
\titleformat{\subsubsection}{\normalsize\bfseries\color{leipGris}}{\thesubsubsection}{0.6em}{}

% --- Portada -----------------------------------------------------------
\newcommand{\logoUniversidad}{img/logo-ucu.png}
\newcommand{\portada}[1]{%
  \begin{titlepage}
  \centering
  \vspace*{1.5cm}
  \IfFileExists{\logoUniversidad}{%
    \includegraphics[width=0.40\textwidth]{\logoUniversidad}\par
    \vspace{1.2cm}%
  }{%
    \fbox{\parbox[c][3cm][c]{0.5\textwidth}{\centering\small\texttt{\logoUniversidad}}}\par
    \vspace{1.2cm}%
  }
  {\Large \proyectoUniversidad\par}
  {\large \proyectoFacultad\par}
  {\large \proyectoCarrera\par}
  \vspace{3cm}
  {\Huge\bfseries\color{leipAzul} #1\par}
  \vspace{0.8cm}
  {\Large \proyectoNombre{} --- \proyectoNombreLargo\par}
  \vspace{2cm}
  {\large Versi\'on \docVersion\par}
  {\large \docFecha\par}
  \vfill
  \begin{tabular}{c}
    Sebastian Forische \\
    Armando Hernandez de la Vega \\
    Luana Acu\~na \\
    Bruno Sebastian Olivera Viera Gomez \\
  \end{tabular}
  \vspace{1.5cm}
  \end{titlepage}
}

% --- Historial de revisiones -------------------------------------------
\newenvironment{historial}{%
  \section*{Historial de revisiones}
  \addcontentsline{toc}{section}{Historial de revisiones}
  \begin{tabular}{@{}l l p{0.52\textwidth} l@{}}
  \toprule
  \textbf{Versi\'on} & \textbf{Fecha} & \textbf{Descripci\'on} & \textbf{Autor} \\
  \midrule
}{%
  \bottomrule
  \end{tabular}
  \bigskip
}

% --- Marcador de contenido pendiente -----------------------------------
\newif\ifpendientes
\pendientestrue

\makeatletter
\newcommand{\listofpendientes}{%
  \section*{Pendientes de este documento}
  \addcontentsline{toc}{section}{Pendientes de este documento}
  \@starttoc{pnd}%
}
\newcommand{\pendiente}[2]{%
  \addcontentsline{pnd}{section}{\protect\numberline{}#1}%
  \ifpendientes
    \begin{tcolorbox}[breakable,colback=leipAmbarBg,colframe=leipAmbar,
                      boxrule=0.6pt,left=6pt,right=6pt,top=4pt,bottom=4pt]
      \textbf{\color{leipAmbar}PENDIENTE.}\ #2
    \end{tcolorbox}
  \fi
}
\makeatother

% =====================================================================
%  ENTORNO PARA DOCUMENTAR UNA DECISION DE ARQUITECTURA
% ---------------------------------------------------------------------
%  Uso:
%
%  \begin{decision}{DA-01}{Titulo de la decision}
%    \dacontexto{Que problema obliga a decidir}
%    \dadecision{Que se decidio}
%    \dajustificacion{Por que}
%    \dasupuestos{En que supuestos se apoya}
%    \darnf{RNF-01, RNF-04}
%    \daalternativas{Que otras opciones se evaluaron y por que se descartaron}
%    \daconsecuencias{Que implica hacia adelante}
%  \end{decision}
% =====================================================================

\newcommand{\dacampo}[2]{%
  \par\noindent
  \begin{tabularx}{\textwidth}{@{}>{\bfseries}p{3.4cm} X@{}}
    #1 & #2
  \end{tabularx}\par\smallskip
}

\newcommand{\dacontexto}[1]{\dacampo{Contexto}{#1}}
\newcommand{\dadecision}[1]{\dacampo{Decisi\'on}{#1}}
\newcommand{\dajustificacion}[1]{\dacampo{Justificaci\'on}{#1}}
\newcommand{\dasupuestos}[1]{\dacampo{Supuestos}{#1}}
\newcommand{\darnf}[1]{\dacampo{Responde a}{#1}}
\newcommand{\daalternativas}[1]{\dacampo{Alternativas}{#1}}
\newcommand{\daconsecuencias}[1]{\dacampo{Consecuencias}{#1}}

\newenvironment{decision}[2]{%
  \subsection*{\color{leipAzul}#1\quad #2}
  \addcontentsline{toc}{subsection}{#1\quad #2}
  \vspace{-0.4em}
  {\color{leipGris}\hrule height 0.5pt}
  \smallskip
}{%
  \bigskip
}

% --- Figura con fallback si la imagen todavia no existe ----------------
\newcommand{\figuraLEIP}[4]{%
  \begin{figure}[htbp]
  \centering
  \IfFileExists{#1}{\includegraphics[width=#4\textwidth]{#1}}{%
    \fbox{\parbox[c][3.5cm][c]{0.8\textwidth}{\centering\color{leipGris}%
      \small\texttt{#1}}}%
  }
  \caption{#2}
  \label{#3}
  \end{figure}
}

\begin{document}

\portada{\docTitulo}

\tableofcontents
\bigskip

\begin{historial}
1.0 & 26/06/2026 & Contenido inicial derivado del anteproyecto & Equipo \\
1.1 & 11/09/2026 & Se completan objetivo, alcance y convenciones, el modelo de datos, el modelo de revisiones, la estrategia de aislamiento, las decisiones de arquitectura y la bibliograf\'ia & Equipo \\
1.2 & 12/09/2026 & Revisi\'on de redacci\'on: los componentes y la infraestructura se describen en presente, igual que las fichas de decisi\'on, y se retiran los p\'arrafos sobre el estado del propio documento & Equipo \\
\end{historial}
\clearpage

% =====================================================================
\section{Introducci\'on}

\subsection{Objetivo del documento}

Este documento describe c\'omo est\'a construido \proyectoNombre{}: qu\'e
componentes lo integran, c\'omo se organizan los datos que procesa, c\'omo se
despliega y, sobre todo, por qu\'e se resolvi\'o de esa manera y no de otra.

Donde la Especificaci\'on de Requerimientos enuncia qu\'e debe hacer el sistema
y con qu\'e restricciones, este documento registra las decisiones t\'ecnicas que
lo hacen posible. Cada decisi\'on relevante queda documentada con su contexto,
las alternativas evaluadas y el requerimiento no funcional al que responde, de
modo que pueda revisarse m\'as adelante sabiendo sobre qu\'e supuestos se tom\'o.

\subsection{Alcance}

El documento cubre la vista general de la arquitectura, los componentes de
\emph{backend}, procesamiento, inteligencia y \emph{frontend}, la organizaci\'on
de los datos y su modelo de revisiones, el uso de modelos de lenguaje dentro de
la ingesta, la seguridad y el aislamiento entre organizaciones, la
infraestructura de despliegue y el registro de decisiones de arquitectura.

No cubre los requerimientos que estas decisiones satisfacen, que se enuncian en
la Especificaci\'on de Requerimientos, ni las interacciones del usuario con el
sistema, que se describen en la Especificaci\'on de Casos de Uso.

\subsection{Convenciones}

Las decisiones de arquitectura se identifican como \textbf{DA-nn} y se registran
en la secci\'on \ref{sec:decisiones}. Los identificadores son estables: una
decisi\'on superada conserva su c\'odigo y se marca como tal, indicando cu\'al
la reemplaza.

Cada decisi\'on indica, en su campo \emph{Responde a}, los requerimientos no
funcionales (\textbf{RNF-nn}) de la Especificaci\'on de Requerimientos que la
motivan.

Las capas de datos se nombran seg\'un el patr\'on Medallion ---Bronze, Silver y
Gold--- complementadas con una zona de cuarentena. El significado de cada una se
define en la secci\'on \ref{sec:datos} y en el Glosario.

% =====================================================================
\section{Vista general de la arquitectura}

\proyectoNombre{} se organiza como un conjunto de componentes desacoplados que
siguen el recorrido de la informaci\'on: un m\'odulo de procesamiento obtiene y
normaliza los datos de las fuentes oficiales, un m\'odulo de inteligencia
calcula sobre ellos, un \emph{backend} centraliza el acceso y los permisos, y un
\emph{frontend} presenta el resultado. Entre ellos, una base relacional y un
almacenamiento de objetos conservan respectivamente los datos estructurados y
los archivos originales.

El diagrama siguiente presenta esos componentes y las relaciones entre ellos,
incluidas las entidades externas con las que el sistema se comunica: las fuentes
oficiales, el proveedor de pagos y los sistemas consumidores.

\figuraLEIP{img/arquitectura.png}{Arquitectura de \proyectoNombre{}}{fig:arquitectura}{0.95}

% =====================================================================
\section{Componentes}

Los componentes se describen a continuaci\'on en el orden en que participan del
recorrido de la informaci\'on. Cada uno se despliega por separado, lo que
permite dimensionarlos y actualizarlos de forma independiente.

\subsection{Backend --- Node.js}

El backend se desarrolla en Node.js y expone una API REST que centraliza la
l\'ogica de negocio, la autenticaci\'on, los permisos por m\'odulo, las
suscripciones y la comunicaci\'on entre el frontend, la base de datos y los
servicios de procesamiento. Expone tambi\'en la API institucional autenticada.

\subsection{M\'odulo de procesamiento de datos --- Python}

El procesamiento de datos se implementa en Python, con Pandas y \emph{jobs}
programados. Esta capa obtiene informaci\'on desde API, Excel, CSV, PDF, HTML o
JSON; valida su contenido, normaliza variables, unidades y fechas, y env\'ia los
registros inconsistentes a cuarentena antes de publicarlos.

\subsection{M\'odulo de inteligencia --- Python}

La capa de inteligencia se implementa en Python mediante un motor
determin\'istico de m\'etricas y un grafo de dependencias para los cruces entre
m\'odulos. Calcula cambios, momentum, divergencias regionales, percentiles
hist\'oricos, sorpresas, persistencia y distancia frente a metas.

El grafo de dependencias cumple una funci\'on adicional a la de ordenar el
c\'alculo: cuando una fuente publica una revisi\'on, determina exactamente
qu\'e m\'etricas y qu\'e cruces quedan desactualizados, de modo que se
recalcule solo lo afectado y no todo el conjunto.

\subsection{Frontend --- React}

El frontend se desarrolla en React, con componentes reutilizables que construyen
\emph{dashboards}, comparadores, fichas por pa\'is, alertas, \emph{briefings} y
paneles administrativos. La interfaz adapta las funcionalidades visibles seg\'un
los permisos y m\'odulos contratados.

La resoluci\'on de permisos es responsabilidad del backend: el frontend adapta
lo que muestra, pero no es la barrera de acceso. Toda consulta se valida
nuevamente del lado del servidor contra el \emph{entitlement} vigente de la
cuenta.

% =====================================================================
\section{Arquitectura de datos}
\label{sec:datos}

\subsection{Patr\'on Medallion}

La informaci\'on se organiza mediante el patr\'on Medallion, que separa el
dato en capas seg\'un su grado de elaboraci\'on. La separaci\'on cumple aqu\'i
una funci\'on concreta: permite reprocesar desde el original cuando se detecta
un error de interpretaci\'on, sin depender de que la fuente todav\'ia publique
ese contenido.

\begin{itemize}
  \item \textbf{Bronze:} conserva los archivos y respuestas originales, inmutables.
  \item \textbf{Silver:} almacena los datos validados, normalizados y trazables.
  \item \textbf{Gold:} contiene las m\'etricas, cruces y se\~nales preparadas para
        \emph{dashboards}, alertas, \emph{briefings}, exportaciones y APIs.
  \item \textbf{Cuarentena:} recibe los registros dudosos antes de que puedan
        publicarse.
\end{itemize}

\subsection{Base de datos --- PostgreSQL}

PostgreSQL almacena los datos normalizados, las m\'etricas, los hist\'oricos, las
revisiones, los usuarios, las organizaciones, los permisos y los registros de
auditor\'ia. Los archivos originales de la capa Bronze se conservan por separado
en almacenamiento de objetos.

\subsection{Almacenamiento de objetos --- Azure Blob Storage}

Azure Blob Storage almacena los archivos originales, los documentos
hist\'oricos, los datos en cuarentena y las exportaciones. As\'i se conserva una
copia verificable de cada fuente sin sobrecargar la base de datos principal.

\subsection{Modelo de datos}

El modelo se organiza en torno a una distinci\'on central: una \emph{serie}
identifica qu\'e se mide ---pa\'is, variable, unidad, frecuencia y horizonte---
mientras que una \emph{observaci\'on} registra un valor concreto de esa serie
para un per\'iodo determinado. Esta separaci\'on es la que permite que dos
pa\'ises sean comparables aunque sus fuentes publiquen con nombres y estructuras
distintas.

La tabla siguiente resume las entidades principales y su papel dentro del
sistema. Cada una se desarrolla en el modelo f\'isico junto con sus claves,
\'indices y restricciones de integridad.

\begin{tabularx}{\textwidth}{@{}l X@{}}
\toprule
\textbf{Entidad} & \textbf{Rol en el sistema} \\
\midrule
Organismo & Instituci\'on emisora: banco central, instituto de estad\'istica o ministerio. \\
\addlinespace
Conjunto de datos & Publicaci\'on concreta de un organismo, con su formato, frecuencia, condiciones de uso y atribuci\'on requerida. Es la unidad sobre la que opera la matriz de fuentes. \\
\addlinespace
Serie & Definici\'on de qu\'e se mide: pa\'is, variable, unidad, frecuencia y horizonte, bajo el modelo com\'un. Una serie puede alimentarse de m\'as de un conjunto de datos a lo largo del tiempo. \\
\addlinespace
Observaci\'on & Valor de una serie para un per\'iodo de referencia, en una versi\'on determinada. Es la unidad m\'inima de dato del sistema. \\
\addlinespace
Ejecuci\'on de ingesta & Registro de una corrida sobre un conjunto de datos: momento, duraci\'on, volumen procesado, \emph{checksum} del original e incidencias detectadas. \\
\addlinespace
Registro en cuarentena & Dato que no super\'o las validaciones, conservado con el motivo del rechazo a la espera de revisi\'on. \\
\addlinespace
M\'etrica & Resultado del c\'alculo determin\'istico sobre una o m\'as series, con sus par\'ametros y las versiones de las observaciones que lo originaron. \\
\addlinespace
Cruce & Relaci\'on calculada entre series de m\'odulos distintos, con las mismas exigencias de trazabilidad que una m\'etrica. \\
\addlinespace
Organizaci\'on & Cuenta corporativa que agrupa usuarios y concentra la suscripci\'on y los permisos. \\
\addlinespace
Usuario & Persona con acceso al sistema, perteneciente o no a una organizaci\'on. \\
\addlinespace
Entitlement & Conjunto de m\'odulos, pa\'ises y funcionalidades habilitados para una cuenta, derivado de su suscripci\'on vigente. \\
\addlinespace
Suscripci\'on & Estado de contrataci\'on de una cuenta, con su identificador externo en el proveedor de pagos y su vigencia. \\
\addlinespace
Registro de auditor\'ia & Asiento de una acci\'on relevante sobre datos, permisos o configuraci\'on, con su actor y su momento. \\
\bottomrule
\end{tabularx}

\pendiente{Completar el modelo f\'isico de datos}{Definir el diagrama entidad-relaci\'on con claves primarias y
for\'aneas, \'indices, restricciones de integridad y estrategia de particionado
para las tablas de observaciones, que son las que crecen con el hist\'orico.}

\subsection{Modelo de revisiones e hist\'oricos}

Las fuentes oficiales revisan sus datos: un valor publicado para un per\'iodo
puede corregirse en publicaciones posteriores, y el valor anterior deja de estar
disponible en el sitio del organismo. Conservar esa historia es una capacidad
central del producto, no un detalle de implementaci\'on, y condiciona c\'omo se
almacenan las observaciones.

El modelo distingue por lo tanto dos ejes temporales independientes para cada
observaci\'on:

\begin{itemize}
  \item \textbf{Per\'iodo de referencia:} el momento al que el dato se refiere;
        por ejemplo, la expectativa de inflaci\'on para diciembre de 2026.
  \item \textbf{Momento de publicaci\'on:} cu\'ando el organismo public\'o ese
        valor, y por lo tanto desde cu\'ando el sistema lo conoce.
\end{itemize}

De esta separaci\'on se derivan tres reglas de almacenamiento:

\begin{itemize}
  \item Las observaciones se escriben solo por agregado. Una revisi\'on no
        modifica el registro anterior: agrega una versi\'on nueva y marca la
        anterior como superada, conservando ambas.
  \item Toda consulta resuelve por defecto la versi\'on vigente de cada
        observaci\'on, pero puede reconstruir el estado de la informaci\'on tal
        como se conoc\'ia en cualquier fecha pasada.
  \item Las m\'etricas registran qu\'e versi\'on de cada observaci\'on
        utilizaron, de modo que un resultado publicado pueda reproducirse
        exactamente aunque las series hayan sido revisadas desde entonces.
\end{itemize}

Esta \'ultima regla es la que hace posible la m\'etrica de sorpresa y el
seguimiento de revisiones: ambas comparan lo que se esperaba en un momento dado
con lo que se sabe despu\'es, y ninguna de las dos puede calcularse si el
sistema conserva solo el \'ultimo valor de cada serie.

% =====================================================================
\section{Uso de modelos de lenguaje}

Los modelos de lenguaje se utilizan en dos puntos acotados del sistema, ambos
con la misma restricci\'on: no producen datos ni conclusiones por s\'i mismos.
El primero interpreta formatos heterog\'eneos que no pueden procesarse de forma
determin\'istica; el segundo redacta, a partir de m\'etricas ya calculadas, el
resumen que acompa\~na al radar. En ning\'un caso un valor num\'erico publicado
por el producto proviene del modelo.

\subsection{Extractor estructurado}

Para resolver la heterogeneidad de formatos entre fuentes, un componente LLM
act\'ua como extractor estructurado: recibe el HTML o archivo publicado
por el banco central, lo limpia y devuelve un JSON normalizado con las
expectativas en el formato interno de \proyectoNombre{}, independientemente de
c\'omo cada banco central publique su informaci\'on.

El extractor se aplica \'unicamente a las fuentes cuya estructura no es estable
o no admite un procesamiento por reglas. Cuando una fuente expone un servicio de
datos o un formato tabular consistente, se procesa de forma determin\'istica y
el modelo no interviene.

\subsection{Generaci\'on del resumen del radar}

Una vez calculadas las m\'etricas, un componente LLM recibe esos resultados ya
procesados y genera autom\'aticamente el resumen descriptivo del radar de
expectativas.

El modelo recibe exclusivamente m\'etricas ya calculadas y validadas, nunca los
datos crudos, y su salida es descriptiva: no introduce cifras propias ni
interpretaciones que no se deriven de las m\'etricas recibidas.

\subsection{Validaci\'on y manejo de errores}

La salida de un modelo de lenguaje puede estar bien formada y ser incorrecta, y
esa diferencia no se detecta por su forma. Por eso ninguna salida del extractor
se publica directamente: atraviesa los mismos controles determin\'isticos que el
resto de la ingesta.

\begin{itemize}
  \item \textbf{Validaci\'on de esquema:} la respuesta debe ajustarse al
        esquema interno esperado en campos, tipos y unidades. Si no lo hace, se
        descarta sin reintentar la interpretaci\'on.
  \item \textbf{Validaci\'on de rango y continuidad:} los valores se contrastan
        con el rango hist\'orico de la serie y con la publicaci\'on anterior.
        Una desviaci\'on fuera de lo esperado env\'ia el registro a cuarentena.
  \item \textbf{Reintentos acotados:} ante una falla de forma se reintenta un
        n\'umero limitado de veces; agotados los reintentos, la ejecuci\'on se
        registra como fallida y se notifica, en lugar de publicar un resultado
        dudoso.
  \item \textbf{Conservaci\'on del original:} el archivo o respuesta de origen
        queda en la capa Bronze, de modo que cualquier extracci\'on pueda
        rehacerse m\'as adelante sin depender de la fuente.
  \item \textbf{Trazabilidad del componente:} cada dato extra\'ido registra que
        su origen fue el extractor, con qu\'e versi\'on del modelo y con qu\'e
        instrucci\'on, para poder reprocesar de forma selectiva ante un cambio.
\end{itemize}

% =====================================================================
\section{Seguridad y aislamiento}

La plataforma expone \'unicamente el frontend y los \emph{endpoints} necesarios.
Incorpora autenticaci\'on, recuperaci\'on de acceso, sesiones o \emph{tokens}
seguros, validaci\'on de solicitudes, gesti\'on de secretos fuera del c\'odigo,
auditor\'ia, roles y permisos por usuario, organizaci\'on, plan, m\'odulo y
pa\'is, y aislamiento multi-\emph{tenant} para las fuentes privadas.

\subsection{Estrategia de aislamiento multi-\emph{tenant}}

El sistema maneja dos clases de informaci\'on con exigencias de aislamiento
distintas, y la estrategia las trata por separado.

La informaci\'on oficial es com\'un a todas las organizaciones: proviene de
fuentes p\'ublicas y no hay raz\'on para duplicarla por cliente. Sobre ella, el
aislamiento es de acceso y no de almacenamiento: las consultas se filtran por el
\emph{entitlement} vigente de la cuenta, resuelto siempre del lado del servidor.

Las fuentes privadas que aporta una organizaci\'on son propias de esa
organizaci\'on y no deben mezclarse con las de ninguna otra. Sobre ellas el
aislamiento es de almacenamiento: cada organizaci\'on tiene su propio espacio
l\'ogico separado, tanto en la base de datos como en el almacenamiento de
objetos, y ning\'un proceso de consulta puede abarcar m\'as de uno.

\begin{itemize}
  \item El identificador de organizaci\'on forma parte de la clave de todo dato
        privado y de la ruta de todo archivo privado, de modo que no exista un
        camino de acceso que lo omita.
  \item Las credenciales de acceso a datos privados se emiten por organizaci\'on
        y no permiten alcanzar el espacio de otra.
  \item Los c\'alculos que combinan datos oficiales con fuentes privadas se
        ejecutan dentro del espacio de la organizaci\'on y su resultado
        permanece all\'i.
  \item Toda consulta sobre datos privados queda asentada en la auditor\'ia con
        su organizaci\'on, su usuario y su alcance.
\end{itemize}

\pendiente{Definir el mecanismo concreto de aislamiento en PostgreSQL}{Decidir entre esquema separado por organizaci\'on, seguridad a
nivel de fila o base separada para las organizaciones que lo requieran por
contrato, evaluando el costo operativo de cada opci\'on frente al n\'umero
esperado de organizaciones. La decisi\'on se registrar\'a como ficha DA en la
secci\'on de decisiones de arquitectura.}

% =====================================================================
\section{Infraestructura y despliegue}

\subsection{Azure Container Apps}

La infraestructura se basa en contenedores Docker desplegados sobre Azure
Container Apps, lo que desacopla los componentes del sistema y permite
desplegarlos y escalarlos por separado.

\subsection{Integraci\'on y despliegue continuo}

La integraci\'on y el despliegue continuo (CI/CD) se implementan con GitHub
Actions.

\subsection{Control de versiones}

El c\'odigo se aloja en GitHub, que provee el control de versiones, la gesti\'on
de incidencias, las revisiones de c\'odigo y el historial completo de cambios.

\subsection{Diagrama de despliegue}

El diagrama de despliegue presenta c\'omo se distribuyen los componentes sobre
la infraestructura: qu\'e contenedores se ejecutan, c\'omo se comunican entre
s\'i y con los servicios gestionados, y qu\'e queda expuesto p\'ublicamente.

\figuraLEIP{img/despliegue.png}{Diagrama de despliegue de \proyectoNombre{}}{fig:despliegue}{0.95}

% =====================================================================
\section{Decisiones de arquitectura}
\label{sec:decisiones}

Esta secci\'on registra las decisiones t\'ecnicas tomadas, con el contexto que
las motiv\'o, las alternativas evaluadas y sus consecuencias. El prop\'osito es
que una decisi\'on pueda revisarse m\'as adelante sabiendo sobre qu\'e supuestos
se tom\'o: si un supuesto deja de ser v\'alido, la ficha indica qu\'e habr\'ia
que reconsiderar.

\begin{decision}{DA-01}{Conservar los originales de forma inmutable}

\dacontexto{Las fuentes oficiales revisan y retiran contenido sin aviso. Si el
sistema conserva \'unicamente el dato ya interpretado, un error de
interpretaci\'on es irrecuperable: el original ya no est\'a disponible en la
fuente para volver a procesarlo.}
\dadecision{Toda respuesta o archivo obtenido de una fuente se conserva sin
modificaciones en la capa Bronze, con su \emph{checksum}, antes de cualquier
procesamiento. Los datos interpretados se derivan siempre de ese original.}
\dajustificacion{Separa el acto de obtener del acto de interpretar. Permite
corregir un extractor y reprocesar todo el hist\'orico sin volver a depender de
la fuente, y permite demostrar de d\'onde sali\'o cada valor publicado.}
\dasupuestos{El volumen de originales conservados es manejable en costo de
almacenamiento, por tratarse de publicaciones peri\'odicas de tama\~no acotado.}
\darnf{RNF-01, RNF-02}
\daalternativas{Conservar solo el dato normalizado: descartada porque hace
irreversible cualquier error de interpretaci\'on. Conservar los originales por
un per\'iodo limitado: descartada porque las revisiones de las fuentes pueden
referirse a per\'iodos antiguos.}
\daconsecuencias{Obliga a un almacenamiento de objetos separado de la base
relacional y a una pol\'itica de nomenclatura y retenci\'on de archivos. A
cambio, habilita el reprocesamiento selectivo como respuesta est\'andar ante
cualquier incidente de calidad de datos.}

\end{decision}

\begin{decision}{DA-02}{Almacenar las observaciones por agregado, no por actualizaci\'on}

\dacontexto{El producto debe poder responder no solo cu\'anto vale hoy una
serie, sino cu\'anto val\'ia seg\'un lo que se sab\'ia en una fecha pasada. Las
m\'etricas de sorpresa y el seguimiento de revisiones dependen de esa
capacidad.}
\dadecision{Una revisi\'on nunca sobrescribe la observaci\'on anterior: agrega
una versi\'on nueva y marca la previa como superada. Cada observaci\'on lleva
per\'iodo de referencia y momento de publicaci\'on como ejes independientes.}
\dajustificacion{Es la \'unica forma de reconstruir el estado de la
informaci\'on en un momento dado y de reproducir exactamente una m\'etrica ya
publicada.}
\dasupuestos{La frecuencia de revisi\'on de las fuentes es moderada, de modo que
el n\'umero de versiones por observaci\'on se mantiene acotado.}
\darnf{RNF-01, RNF-02}
\daalternativas{Actualizar el valor en su lugar: descartada porque elimina la
historia que el producto necesita exponer. Registrar solo la \'ultima revisi\'on
anterior: descartada porque no permite reconstruir series con m\'as de una
correcci\'on.}
\daconsecuencias{Las tablas de observaciones crecen con el n\'umero de
revisiones y requieren \'indices y particionado adecuados. Las consultas deben
resolver expl\'icitamente qu\'e versi\'on piden, con la vigente como
comportamiento por defecto.}

\end{decision}

\begin{decision}{DA-03}{Resolver los permisos en el servidor mediante \emph{entitlements}}

\dacontexto{El acceso a m\'odulos, pa\'ises y funcionalidades depende de lo que
cada cuenta tenga contratado, y ese alcance cambia con cada contrataci\'on,
renovaci\'on o cancelaci\'on. Adem\'as, el mismo contenido puede quedar
bloqueado por una restricci\'on regulatoria ajena a la contrataci\'on.}
\dadecision{El backend resuelve, para cada solicitud, el \emph{entitlement}
vigente de la cuenta y filtra el resultado antes de devolverlo. El frontend
adapta lo que muestra, pero no constituye la barrera de acceso.}
\dajustificacion{Concentra la decisi\'on de acceso en un \'unico lugar, aplicable
por igual a la interfaz web y a la integraci\'on program\'atica, y permite que
un bloqueo regulatorio tenga efecto inmediato sobre todas las v\'ias de acceso.}
\dasupuestos{La resoluci\'on del \emph{entitlement} es suficientemente
econ\'omica como para ejecutarse en cada solicitud o para cachearse por
sesi\'on sin riesgo de quedar obsoleta.}
\darnf{RNF-04, RNF-07}
\daalternativas{Filtrar en el cliente: descartada por no ser una barrera real.
Precalcular vistas por plan: descartada porque multiplica el almacenamiento y
retrasa el efecto de los bloqueos.}
\daconsecuencias{Todo nuevo camino de acceso a datos debe pasar por la misma
resoluci\'on de permisos. Los cambios de contrataci\'on y los bloqueos se
reflejan sin intervenci\'on adicional.}

\end{decision}

\begin{decision}{DA-04}{Delegar el cobro en un proveedor externo con \emph{checkout} alojado}

\dacontexto{El producto necesita cobrar suscripciones, pero almacenar o procesar
datos de tarjetas impondr\'ia exigencias de cumplimiento desproporcionadas
respecto del volumen esperado.}
\dadecision{El pago se resuelve \'integramente en el \emph{checkout} alojado del
proveedor. El sistema conserva solo el identificador externo, el estado de la
suscripci\'on y los datos necesarios para actualizar los \emph{entitlements}, y
recibe los cambios de estado por \emph{webhook}.}
\dajustificacion{Elimina los datos de tarjeta del alcance del sistema y reduce
la superficie de cumplimiento al m\'inimo.}
\dasupuestos{Existe un proveedor que ofrece \emph{checkout} alojado y
notificaciones por \emph{webhook} en las plazas de inter\'es.}
\darnf{RNF-08}
\daalternativas{Integrar el cobro dentro de la aplicaci\'on: descartada por las
exigencias de cumplimiento que traer\'ia. Cobro fuera de l\'inea: descartada por
no sostener la renovaci\'on autom\'atica.}
\daconsecuencias{Introduce una dependencia externa en el camino de
contrataci\'on, registrada como riesgo R-06. Obliga a validar cada
notificaci\'on recibida contra la intenci\'on registrada y a prever la
reconciliaci\'on de estados.}

\end{decision}

\begin{decision}{DA-05}{Usar modelos de lenguaje solo para interpretar y describir, nunca para calcular}

\dacontexto{Las fuentes publican en formatos que no siempre admiten
procesamiento por reglas, y un modelo de lenguaje resuelve bien esa
interpretaci\'on. Pero su salida puede ser incorrecta sin dejar se\~nales de
serlo, y el producto entrega informaci\'on que el usuario cita.}
\dadecision{El modelo se usa en dos puntos acotados: interpretar fuentes no
estructuradas y redactar el resumen del radar a partir de m\'etricas ya
calculadas. Ning\'un valor num\'erico publicado proviene del modelo, y toda
salida atraviesa validaciones determin\'isticas antes de publicarse.}
\dajustificacion{Aprovecha la capacidad del modelo donde aporta ---la
heterogeneidad de formatos--- sin trasladarle la responsabilidad sobre las
cifras, que es la base de la confianza en el producto.}
\dasupuestos{Las validaciones determin\'isticas de esquema, rango y continuidad
son suficientes para detectar las salidas incorrectas m\'as probables.}
\darnf{RNF-01, RNF-10}
\daalternativas{Procesar todo por reglas: descartada por el costo de mantener un
extractor a medida por fuente inestable. Permitir que el modelo calcule o
interprete m\'etricas: descartada por no poder garantizar reproducibilidad.}
\daconsecuencias{Obliga a mantener el conjunto de validaciones como parte
central de la ingesta y a registrar versi\'on de modelo e instrucci\'on junto a
cada dato extra\'ido, para poder reprocesar de forma selectiva.}

\end{decision}

\begin{decision}{DA-06}{Desplegar los componentes en contenedores desacoplados}

\dacontexto{Los componentes tienen perfiles de carga distintos: la ingesta se
ejecuta por r\'afagas seg\'un el calendario de publicaci\'on de las fuentes,
mientras que la aplicaci\'on web atiende consultas de forma continua.}
\dadecision{Cada componente se empaqueta como contenedor y se despliega por
separado sobre Azure Container Apps, con integraci\'on y despliegue continuo
mediante GitHub Actions.}
\dajustificacion{Permite dimensionar y actualizar cada componente de forma
independiente, y evita que una ejecuci\'on de ingesta pesada degrade la
respuesta de la aplicaci\'on web.}
\dasupuestos{La infraestructura del proveedor se mantiene disponible dentro de
las condiciones previstas.}
\darnf{RNF-04}
\daalternativas{Despliegue monol\'itico: descartada por acoplar los perfiles de
carga. Funciones sin servidor para la ingesta: descartada por los l\'imites de
tiempo de ejecuci\'on frente a procesamientos largos.}
\daconsecuencias{Introduce la dependencia del proveedor de nube, registrada como
riesgo R-07, y exige definir la comunicaci\'on y el manejo de errores entre
componentes desplegados por separado.}

\end{decision}

\pendiente{Registrar las decisiones de arquitectura pendientes}{Quedan por documentar como fichas DA: el mecanismo concreto de
aislamiento multi-\emph{tenant} en PostgreSQL, la estrategia de particionado de
las tablas de observaciones, el mecanismo de programaci\'on de las ejecuciones
de ingesta y el criterio de invalidaci\'on de cach\'e ante revisiones.}

% =====================================================================
\section{Dise\~no}

Esta secci\'on cubre el dise\~no interno de cada componente: la estructura de
m\'odulos del backend, la organizaci\'on del motor de m\'etricas y su grafo de
dependencias, la composici\'on del frontend y los contratos de comunicaci\'on
entre componentes.

\pendiente{Desarrollar el dise\~no interno de los componentes}{Documentar la estructura de m\'odulos del backend y sus contratos,
la organizaci\'on del motor determin\'istico de m\'etricas y del grafo de
dependencias, la composici\'on de componentes del frontend y el manejo de
errores entre componentes desplegados por separado.}

% =====================================================================
\section{Bibliograf\'ia}

Las referencias siguientes sustentan las decisiones registradas en este
documento. Para cada una se indica para qu\'e se utiliz\'o.

\begin{tabularx}{\textwidth}{@{}X X@{}}
\toprule
\textbf{Referencia} & \textbf{Para qu\'e se utiliza} \\
\midrule
Michael Nygard, \emph{Documenting Architecture Decisions} (2011). &
Formato y criterio de las fichas de decisi\'on de la secci\'on
\ref{sec:decisiones}: registrar contexto, decisi\'on, alternativas y
consecuencias en lugar de solo el resultado. \\
\addlinespace
Documentaci\'on de la arquitectura Medallion, Databricks. &
Definici\'on de las capas Bronze, Silver y Gold y del criterio de separaci\'on
por grado de elaboraci\'on del dato. \\
\addlinespace
Documentaci\'on oficial de PostgreSQL. &
Capacidades de indexado, particionado y control de integridad consideradas para
el modelo de observaciones y revisiones. \\
\addlinespace
Documentaci\'on oficial de Azure Container Apps y Azure Blob Storage. &
Modelo de despliegue en contenedores y caracter\'isticas del almacenamiento de
objetos utilizado para la capa Bronze. \\
\addlinespace
Documentaci\'on oficial de Node.js, React y Pandas. &
Capacidades de las tecnolog\'ias seleccionadas para \emph{backend},
\emph{frontend} y procesamiento de datos. \\
\addlinespace
Documentaci\'on oficial de GitHub Actions. &
Definici\'on del proceso de integraci\'on y despliegue continuo. \\
\bottomrule
\end{tabularx}

\pendiente{Completar la bibliograf\'ia con URL y fechas de consulta}{Registrar para cada referencia su URL y su fecha de consulta,
siguiendo el mismo criterio que la matriz de fuentes aplica a los conjuntos de
datos.}

% =====================================================================
\clearpage
\listofpendientes

\end{document}
